\origpage{325}{N. Robertson, P.D. Seymour / Journal of Combinatorial Theory, Series B 90 (2004) 325--385}
{\centering
{\Large Graph Minors. XIX. Well-quasi-ordering on a surface\par}\medskip

Neil Robertson$^{a,1}$ and P.D. Seymour$^{b,2}$\par
$^a$\textit{Department of Mathematics, Ohio State University, 231 West 18th Ave., Columbus, OH 43210, USA}\par
$^b$\textit{Telcordia Technologies, 445 South St., Morristown, NJ 07960, USA}\par
Received 13 February 1996\par
}
\begin{quote}
\textbf{Abstract}

In a previous paper (\textit{J. Combin. Theory} 48 (1990) 255) we showed that for any infinite set of (finite) graphs drawn in a fixed surface, one of the graphs is isomorphic to a minor of another. In this paper we extend that result in two ways:

\hspace*{12pt}$\bullet$ we generalize from graphs to hypergraphs drawn in a fixed surface, in which each edge has two or three ends, and

\hspace*{12pt}$\bullet$ the edges of our hypergraphs are labelled from a well-quasi-order, and the minor relation is required to respect this order.

This result is another step in the proof of Wagner's conjecture, that for any infinite set of graphs, one is isomorphic to a minor of another.
\end{quote}

\sectionhead{1. Introduction}

This paper is the penultimate step in the proof of Wagner's conjecture, that for any infinite set of finite graphs, one is isomorphic to a minor of another [4]. Roughly speaking (we shall define our terms later), we wish to show that for any fixed surface $\Sigma$ and fixed well-quasi-order $\Omega$, if $H_1, H_2, \dots$ is an infinite sequence of hypergraphs drawn in $\Sigma$ where each edge has two or three ends, and for each $i \ge 1$, $\phi_i \colon E(H_i) \to E(\Omega)$ is some function, then there exist $j > i \ge 1$ such that $H_i$ is isomorphic to a ``minor'' of $H_j$, and for each edge $e$ of $H_i$ the corresponding edge $f$ of $H_j$ satisfies $\phi_i(e) \le \phi_j(f)$. There are also special rules concerning the boundary of $\Sigma$ and edges%
\origfootnote[1]{This research was partially supported by NSF Grant DMS 8504054.}%
\origfootnote[2]{Present address: Mathematics Department, Princeton University, Princeton, NJ 08544, USA.}
